%% ============================================================
%% sections/01_introduction.tex
%%
%% The landscape, the obstruction quoted in the authors' own words, the
%% main theorem, the architecture of its proof, and the provenance
%% paragraph. Target: three pages.
%% ============================================================

\section{Introduction}
\label{sec:intro}

\subsection{The landscape}
\label{sec:landscape}

Let $p$ be a prime, $q$ a power of $p$, and $X/\Fq$ a smooth affine curve
with smooth compactification $\Xbar$ of genus $g$ and $\Sinf = \Xbar
\setminus X$. For a nontrivial finite character $\rho$ of $\pi_1(X)$, the
$L$-function $L(\rho,s)$ is a polynomial, and the Newton-over-Hodge problem
asks for a lower bound on its $q$-adic Newton polygon $\NP_q(\rho)$ in terms
of the ramification of $\rho$ alone.

Over the projective line and its affinoids the problem is closed at every
prime, $p = 2$ included. Zhu's theorem on one-dimensional
affinoids~\cite{zhu2004newton} is sharp for any prime coprime to the pole
orders; Liu and Wan's $T$-adic bound~\cite[Theorem~5.2]{liuwan2009tadic}
carries no hypothesis on $p$ whatsoever, and neither does Schmidt's
affinoid version~\cite{schmidt2023tadic}; the Davis--Wan--Xiao
tower~\cite{daviswanxiao2016newton} needs only $\gcd(d,p) = 1$. Nothing
below is a claim about $\PP^1$. Behind these sit the Newton-polygon theory
of exponential sums of Adolphson and Sperber~\cite{adolphsonsperber1989exponential,adolphsonsperber1993twisted}
and Wan~\cite{wan1993newton,wan2004variation}, and the stratification
results of Blache and F\'erard~\cite{blacheferard2007newton}; the $T$-adic
line continues in~\cite{haessig2026}.

Over an \emph{arbitrary} smooth affine curve the situation is different,
and the difference is the point of this paper. The framework that produces
a \emph{ramification-defined} Hodge bound there --- one built from the Swan
conductors of $\rho$ and nothing else --- is Kramer-Miller's and
Kramer-Miller--Upton's. It comprises~\cite{kramermiller2021padic} for exponential sums,
\cite{kramermiller2020abelian} for abelian Artin $L$-functions, and~\cite{kramermillerupton2021newton,kramermillerupton2021variation} for the local-to-global
contact theory and its variation in families. Every one of these papers
assumes $p$ odd. The hypothesis is printed in the first section of each:
\emph{``Let $p$ be a prime with $p \ge 3$''}~\cite{kramermiller2020abelian}, \emph{``Let $p$ be an odd prime and let $q$
be a power of $p$''}~\cite{kramermillerupton2021newton}.

The exclusion is neither an oversight nor an unattempted case. The authors
name its two components.

\subsection{The obstruction, in the authors' words}
\label{sec:obstruction}

The first component is analytic. At a point $P$ where the auxiliary Belyi
map $\eta$ takes the value $1$, the relevant Frobenius is not a $p$-power
map, and the Dwork operator $\Up$ moves poles around in a way that has to
be controlled by a weight on a growth module. Kramer-Miller and Upton
carried out the $p = 2$ computation and recorded the outcome in their
Remark~6.5, which we quote in full because it is the statement this paper
repairs.

\begin{quotation}
\noindent
Suppose that $p = 2$. For $k \ge 3$, define $\wt(k) = \lfloor (k-1)/3
\rfloor$. A similar construction provides a submodule $\Am \subset
A^{\dagger}_{\pi,P}$ with the following property: Let $k = 2\ell - r$ with
$r = 0$ or $1$. Then
\[
  \Up\bigl(\pi^{\wt(k)/m_P} t_P^{-k}\bigr) \in
  \pi^{(\wt(k) - \wt(\ell + r))/m_P}\, \Am .
\]
\emph{This estimate is too low for applications to the global setting.} For
example, if $k = 5 = 2\cdot 3 - 1$, then $\wt(k) - \wt(\ell+r) = 0$, and
this contributes an extra segment of slope $0$ in the global Hodge bound
below.
\hfill~\cite[Remark~6.5]{kramermillerupton2021newton}
\end{quotation}

The estimate is not merely weak: it breaks a \emph{count}. The
Deuring--Shafarevich argument that supplies the slope-$0$ part of the
polygon pins the number of slope-$0$ segments exactly, and a single index
with zero gain inserts one more. That is why the remark says ``too low for
applications'' rather than ``weaker than we would like''.

The second component is geometric, and was stated two years earlier:

\begin{quotation}
\noindent
Finally, we mention our requirement that $p \ge 3$. When $p = 2$ it is
likely that the methods in this paper still work. The main difficulty is
that some estimates in section~4 must be modified. It is also not
immediately clear that we can find a cover $\eta : X \to \PP^1_{\Fq}$
satisfying the desired properties. To construct $\eta$, we use the fact
that $X$ admits a simply branched map to $\PP^1_{\Fq}$, which is false when
$p = 2$.
\hfill~\cite[\S1.4]{kramermiller2021padic}
\end{quotation}

The construction that fails is a composition ending in the $(p-1)$-power
map $z \mapsto z^{p-1}$, whose job is to create a fibre of uniform
ramification index $p-1$ over the point $1$. At $p = 2$ that map is the
identity, and the whole stage evaporates.

The gap is live rather than historical. Booher, Groen and
Kramer-Miller~\cite{boohergroenkramermiller2025} treat $\Z/2$-covers in
characteristic two through their Ekedahl--Oort types rather than the full
$2$-adic Newton polygon, which is the move one makes when the $2$-adic
polygon is out of reach.

\subsection{The result}
\label{sec:result}

The obstruction of Remark~6.5 is not analytic. The number $3$ in $\lfloor
(k-1)/3 \rfloor$ is the tame ramification index of the auxiliary map, and
the estimate is repaired by a different weight on the same module ---
Kramer-Miller--Upton's own weight plus a parity indicator,
\[
  \wt(k) = \Bigl\lfloor \frac{k-1}{3} \Bigr\rfloor + (k \bmod 2),
\]
which we prove admissible for every $k$, with defect $\dfct(k) \ge 1$ and
$\dfct(k) \sim k/6$, that rate being exactly optimal
(\S\ref{sec:local}--\S\ref{sec:sharpness}). Together with a
characteristic-$2$ tame Belyi map of uniform index $3$ over the point $1$
(\S\ref{sec:tame-map}), this yields the following.

\begin{theorem}[Main theorem]
\label{thm:main}
Let $q = 2^a$, let $X/\Fq$ be a smooth affine curve with smooth
compactification $\Xbar$ of genus $g$ and $\Sinf = \Xbar \setminus X$, and
let $\rho$ be a nontrivial finite character of $\pi_1(X)$ of any $2$-power
order $2^n$. Then
\[
  \NP_q(\rho) \ \succeq \ \HP_q(\rho)
\]
as full polygons: with the ramification-defined Hodge polygon of
\textup{\cite[\S1.2]{kramermillerupton2021newton}}, no truncation, and no
restriction on $n$.
\end{theorem}

Here $\HP_q(\rho)$ is the Kramer-Miller ramification-defined Hodge polygon,
the polygon with slope multiset
\begin{equation}
\label{eq:hodge-slopes}
  \{0\}^{\,g-1+|\Sinf|} \ \cup\ \{1\}^{\,g-1+|\Sinf|} \ \cup\
  \bigcup_{P \in \Sinf}
  \Bigl\{ \tfrac{1}{\dP},\ \tfrac{2}{\dP},\ \dots,\ \tfrac{\dP-1}{\dP} \Bigr\},
\end{equation}
where $\dP$ is the Swan conductor of $\rho$ at $P$ and $|\Sinf|$ is counted
geometrically; $\succeq$ is the relation ``lies on or above'', as in the
sources. \Cref{thm:main} is the exact characteristic-$2$ analogue of~\cite[Theorem~1.1]{kramermiller2020abelian} for characters of $p$-power
order, with the same polygon and the same generality of base.

The local-to-global \emph{contact} criterion of~\cite[Theorem~1.1]{kramermillerupton2021newton} --- the two polygons touch globally
if and only if they touch locally at every wild point --- also goes
through, but only on an initial segment.

\begin{theorem}[Contact, on an initial segment]
\label{thm:contact}
Assume $X$ ordinary and let $\rho$ have order $2^n$. For every $q$-adic $r$
with $0 \le r \le 2^{1-n}$, the polygons $\HP_q^{<r}(\rho)$ and
$\NP_q^{<r}(\rho)$ have the same terminal point if and only if
$\HP_q^{<r}(\rhoP)$ and $\NP_q^{<r}(\rhoP)$ have the same terminal point
for every $P \in \Sinf$. In particular the full range $r \in [0,1]$ of
\textup{\cite[Theorem~1.1]{kramermillerupton2021newton}} is covered at order $2$.
\end{theorem}

The restriction in \Cref{thm:contact} is not an artifact of our route. We
prove in \S\ref{sec:sharpness} that the index $k = 2e$ is a self-loop of the
type-$2$ operator with transition coefficient exactly $\SelfLoopCoeff$, for
\emph{every} odd tame index $e$ and independently of $e$; so $\dfct(2e) \le
1$ for every weight whatsoever, and no choice of auxiliary index relieves
the cap. \Cref{thm:contact} rests in addition on an exactness statement that~\cite{kramermillerupton2021newton} asserts without proof, at every $p$ and for its own
weight; \Cref{thm:main} does not (\S\ref{sec:assembly}).

\subsection{Architecture of the proof}
\label{sec:architecture}

The global machinery is entirely Kramer-Miller's. What this paper supplies
is the pair of ingredients their own papers name as missing, and the
verification that nothing else in the argument sees the prime.

\begin{enumerate}[label=(\arabic*),leftmargin=*,itemsep=2pt]
\item \textbf{The type-$2$ operator, solved (\S\ref{sec:local}).} Writing
  $\Utwo(t^{-k}) = \sum_j c_{k,j}t^{-j}$ at odd tame index $e$, the support
  is a single arithmetic progression $\jm(k) + e\Z_{\ge 0}$ and the
  coefficients along it are the Taylor coefficients of $\cosh(\lambda
  \arcsinh z)$ and $\sinh(\lambda\arcsinh z)/z$ with $\lambda = k/e$
  (\Cref{thm:closed-form}); the proof is a change of variable turning the
  defining trace identity into one second-order ODE\@.
\item \textbf{Their valuations, exactly.} Legendre's formula converts the
  closed form into an exact expression for $\vtwo(c_{k,j})$
  (\Cref{thm:valuation}), whence a four-case parity argument gives the tail
  estimate $\vtwo \ge m$, tight only at $k \equiv 2 \bmod 4$, $m = 1$
  (\Cref{lem:tail}).
\item \textbf{The repaired weight.} That tightness and the failure of
  $\lfloor (k-1)/3 \rfloor$ occur at the same indices and in opposite
  directions, which is why a parity indicator --- and nothing more --- fixes
  the estimate (\Cref{thm:weight}).
\item \textbf{Sharpness (\S\ref{sec:sharpness}).} One coefficient,
  $c_{2e,2e} = \SelfLoopCoeff$, caps every weight at rate $\RateGeneral$;
  at $e = 3$ that is $\RateGamma$, which the repaired weight attains.
\item \textbf{The geometric input (\S\ref{sec:tame-map}).} Replacing the
  $(p-1)$-power stage by an $e$-power stage, and Fulton's simply-branched
  map by the tame map of Kedlaya--Litt--Witaszek and Sugiyama--Yasuda,
  produces the required Belyi map at $p = 2$ (\Cref{lem:belyi}).
\item \textbf{Assembly (\S\ref{sec:assembly}).} The two ingredients go into
  \cite{kramermiller2020abelian}, whose remaining hypotheses are
  $p$-uniform; the weight transports through a dictionary between the two
  papers' growth modules that holds on the nose.
\end{enumerate}

En route we record four corrections to the sources, none of which is a $p =
2$ obstruction and none of which invalidates a published theorem: a
paraphrase in~\cite{kramermillerupton2021newton} that over-states the freedom in
choosing a linear map (\S\ref{sec:c-correction}), a strict inequality in~\cite[Definition~6.3]{kramermillerupton2021newton} inconsistent with its own claim,
and two statements in~\cite[\S4.1]{kramermiller2020abelian} that are false
as printed and true in a repaired form (\S\ref{sec:defects}). All four are
$p$-uniform.

\subsection{Provenance, machine verification, and what is not claimed}
\label{sec:provenance}

This paper was produced by a multi-agent artificial-intelligence workflow
under human direction. Separate agents extracted and quoted the sources,
attacked the problem, adversarially audited the results, and drafted the
exposition; the audit was run as a hostile re-derivation rather than a
review, and it refuted or downgraded several intermediate claims, including
one range statement that had already propagated into a dependent draft. Every
theorem stated here was re-derived independently of the draft that first
proposed it, and every computational assertion is backed by at least two
implementations sharing no code, in exact rational arithmetic. The
replication package carries these programs --- each asserts its findings and
exits nonzero on failure --- together with the complete research log and an
error ledger recording what the project got wrong before it got it right.

\emph{The proofs in this paper have not yet been read by a human referee.}
We state that once, here, rather than qualifying each theorem: the
distinction that matters to a reader is not which statement carries which
internal label, but that the whole has had machine verification and
adversarial re-derivation and has not had peer review. Every quotation from
a source was transcribed from the published or arXiv text by at least two
independent readers; the passage in \S\ref{sec:obstruction} was transcribed
five times.

We claim nothing for Newton over Hodge at $p = 2$ over $\PP^1$ or over
affinoids, which is published (\S\ref{sec:landscape}). We do not re-verify
Deuring--Shafarevich, Katz--Gabber, Monsky's trace formula, the functional
analysis of~\cite[\S6]{kramermiller2020abelian}, or the internal proofs of~\cite{kedlayalittwitaszek2020tame} and~\cite{sugiyamayasuda2020belyi}; these
are cited as their users cite them. Section~\ref{sec:remarks} states what
remains open.

\subsection{Organisation}

\Cref{sec:local} solves the type-$2$ operator and proves the repaired weight
admissible. \Cref{sec:sharpness} shows the rate is exactly $\RateGamma$ and
draws the consequence for the contact theory. \Cref{sec:tame-map}
constructs the characteristic-$2$ Belyi map and settles the arithmetic
condition it needs. \Cref{sec:assembly} assembles \Cref{thm:main} and
records the source defects met on the way. \Cref{sec:remarks} states the
open problems.
